\documentclass[11pt]{article}

\usepackage[margin=25mm]{geometry}

\usepackage{amsmath,amssymb,fontspec,unicode-math}

\IfFontExistsTF{DejaVu Serif}{\setmainfont{DejaVu Serif}}{\setmainfont{Latin Modern Roman}}

\IfFontExistsTF{Latin Modern Math}{\setmathfont{Latin Modern Math}}{\setmathfont{DejaVu Math TeX Gyre}}

\usepackage{xurl,hyperref}

\setlength{\parindent}{0pt}\setlength{\parskip}{6pt}

\begin{document}

\title{Alligat0R: Pre-Training through Covisibility Segmentation for Relative Camera Pose Regression}

\author{ELVIA preliminary AI analysis}\date{21 October 2026}\maketitle

Preliminary assessment: No Issue Found.\par

Independent scientific review remains pending. Author adoption has not been requested.

No substantive scientific error was established in the inspected covisibility objective, pose losses, metric interpretation or reported split descriptions. A released host timer remains a limited reproducibility caveat without a replacement runtime measurement.

\section*{Review scope}

Claim-focused reading of the original paper and relevant author-released passages; direct mathematical or static-source reasoning. No author program, earlier verifier, model or benchmark was executed during this check.

\section*{Problem 1: The pose-loss normalization admits an equivalent parameterization}

Source: Pose loss, arXiv v2 Eq.(8), PDF p.5

Assessment: not\_issue; impact: minor

\subsection*{Source claim (paraphrase)}

The manuscript uses squared residual norm with uncertainty, whereas code can use component-mean MSE.

\subsection*{Analysis}

The code form \(e^{-s}\|r\|^2/d+s/2\) equals \(\|r\|^2/(2\sigma^2)+\log\sigma+\tfrac12\log(2/d)\) under \(s=\log(2\sigma^2/d)\). The added constant does not alter optimization. A sign-folded angular diagnostic also does not replace metric translation error: \((1,0,0)\) versus \((-1,0,0)\) has distance \(2\).

\subsection*{Finding and scope}

These are compatible parameterization/metric distinctions, without a demonstrated scientific contradiction.

\section*{Problem 2: Unsynchronized host timing is only a limited caveat}

Source: Released eval.py, get\_pose\_estimation, lines 59–62; runtime comparison

Assessment: not\_issue; impact: minor

\subsection*{Source claim (paraphrase)}

The release measures host elapsed time around a potentially asynchronous CUDA forward call.

\subsection*{Analysis}

The timer stops before outputs are copied to CPU, so it need not include all completed GPU work. However historical timing provenance, device execution and the complete measurement path are unestablished. No numerical runtime was remeasured.

\subsection*{Finding and scope}

A synchronized/event-based rerun would be needed to assess the published runtime; the current observation alone does not refute it.

\section*{Limits}

Preliminary AI assessment; independent human scientific adjudication remains pending.\par

Analytic witnesses and static release checks do not remeasure the published benchmark results.\par

Scene memberships, raw annotations, trained-model accuracy and GPU runtime were not independently reproduced.\par

\section*{Source identity}

arXiv:2503.07561v2, 1 December 2025

\url{https://arxiv.org/pdf/2503.07561v2}

Reviewed PDF SHA-256: \nolinkurl{c63b9b9ed0bb1b980ab9d9dd3f06e9f3756610db214c6d525032b33bf56c86f6}

The SHA-256 identifies the exact retained PDF inspected in this rapid source/argument check. Original supplied report bytes are retained locally. No unambiguous original PDF retrieval timestamp was recovered; file modification time was not substituted.

\end{document}
